\section{引言：从 RNN 到 Transformer}

本讲将介绍深度学习中最重要的两个概念：\textbf{注意力机制}（Attention）和\textbf{Transformer 架构}。Attention 是一种全新的神经网络计算原语，本质上操作于\textbf{向量集合}之上；Transformer 则是以 self-attention 为核心的神经网络架构。

\begin{importantbox}{Transformer 的统治地位}
今天几乎所有大规模的深度学习应用——图像分类、图像生成、文本生成、文本分类、音频处理——都在使用 Transformer 架构。它自 2017 年提出以来，已经跨越了多个数量级的规模增长（从 2 亿参数到数万亿参数），且架构本身变化甚微。
\end{importantbox}

尽管 Transformer 是现代深度学习的基石，但注意力机制的思想\textbf{源于 RNN}。我们将沿着历史脉络，从 RNN 的瓶颈问题出发，逐步推导出注意力机制和 Transformer 架构。

\begin{figure}[H]
\centering
\includegraphics[width=0.95\textwidth]{slides-images/slide-001.jpg}
\caption{课程路线图：从 RNN 的瓶颈问题出发，引入注意力机制，最终构建 Transformer 架构\protect\footnotemark}
\end{figure}
\footnotetext{Slides 第2页。}

\subsection{本章小结}

注意力机制和 Transformer 是当代深度学习的核心。理解它们不仅需要掌握数学公式，更要理解其设计动机——如何克服 RNN 的固有局限。

%% ========== Seq2Seq 与瓶颈 ==========

